\section{Introduction}
%

%
%

Bilinear maps stand at the forefront of many basic questions in combinatorics and theoretical computer science. 
A bilinear map is, intuitively, just a collection of matrices. 
Formally, a bilinear map $f \colon \FF^{n_1} \times \FF^{n_2} \to \FF^m$, where $\FF$ is any field, is a map $f(\B{x},\B{y})=(f_1(\B{x},\B{y}),\ldots,f_m(\B{x},\B{y}))$ whose every component $f_k$ is a bilinear form $f_k(\B{x},\B{y}) = \sum_{i,j} a_{i,j,k}x_iy_j$, or equivalently, $\B{x}^T A_k \B{y}$ for some matrix $A_k \in \FF^{n_1 \times n_2}$.
While linear maps are thoroughly understood thanks to linear algebra, 
bilinear maps are---in more than one way---still very much a mystery.



\subsection{Structure vs.\ randomness}

In this paper we prove a tight relation between the \emph{slice rank} and the \emph{analytic rank} of bilinear maps, or \emph{$3$-tensors}.
Our proof crucially uses the notion of \emph{geometric rank} as an intermediary, enabling the use of tools from algebraic geometry to ultimately prove that these three notions of rank are in fact equivalent up to a constant.



A $3$-tensor (or sometimes simply a tensor) over a field $\FF$ is a three-dimensional matrix $(a_{i,j,k})_{i,j,k}\in\FF^{n_1 \times n_2 \times n_3}$ with entries $a_{i,j,k} \in \FF$. 
Equivalently, a tensor can be thought of as a degree-$3$ polynomial, namely, a trilinear form $T(\B{x},\B{y},\B{z}) = \sum_{i,j,k} a_{i,j,k} x_iy_jz_k$ with coefficients $a_{i,j,k} \in \FF$, where $\B{x}=(x_1,\ldots,x_{n_1})$, $\B{y}=(y_1,\ldots,y_{n_2})$, $\B{z}=(z_1,\ldots,z_{n_3})$.\footnote{A trilinear form means that every monomial has exactly one variable from $\B{x}$, one from $\B{y}$, and one from $\B{z}$, and so is linear separately in each of $\B{x}$,$\B{y}$, and $\B{z}$.}
Note that a tensor is just a symmetric way to think of a bilinear map 
$f \colon \FF^{n_1} \times \FF^{n_2} \to \FF^{n_3}$, where $f=(f_1,\ldots,f_{n_3})$ with $f_k(\B{x},\B{y}) = \sum_{i,j} a_{i,j,k}x_iy_j$; indeed, each $f_k$ corresponds to a slice $(a_{i,j,k})_{i,j}$ of $T$.\footnote{Yet another point of view is that a $3$-tensor is a member of the vector space $V_1 \otimes V_2 \otimes V_3$, where $V_1=\FF^{n_1},V_2\in\FF^{n_2},V_3=\FF^{n_3}$ are finite-dimensional vector spaces over $\FF$.}
As opposed to matrices, which have only one notion of rank, there are multiple notions of rank for $3$-tensors. 
The notions of rank of $3$-tensors we consider are defined as follows:
\begin{itemize}
	\item The slice rank of $T$, denoted $\SR(T)$, is the smallest $r \in \NN$ such that $T$ can be decomposed as $T = \sum_{i=1}^r f_i g_i$ where $f_i$ is an $\FF$-linear form in either the $\B{x}$, $\B{y}$, or $\B{z}$ variables and $g_i$ is an $\FF$-bilinear form in the remaining two sets of variables, for each $i$.
	\item The analytic rank of $T$ over a finite field $\FF$ is given by $\AR(T) = -\log_{|\FF|}\Exp_{\B{x},\B{y},\B{z}} \chi(T(\B{x},\B{y},\B{z}))$,\footnote{We use $\Exp_\B{x}$ to denote averaging, so $\Exp_\B{x \in \FF^{n}}$ stands for $|\FF|^{-n}\sum_{\B{x} \in \FF^{n}}$.} where we fix $\chi$ to be any nontrivial additive character of $\FF$ (e.g., $\chi(x)=\exp(2\pi i x/p)$ when $\FF=\FF_p$ is prime). 
	\item The geometric rank of $T$, viewed as a bilinear map $f$,\footnote{Although permuting $\B{x},\B{y},\B{z}$ gives rise to three distinct bilinear maps corresponding to $T$, the definition of $\GR(T)$ is invariant under them (see Theorem~3.1 in~\cite{KoppartyMoZu20}).} is defined as $\GR(T) = \codim \ker f$, the codimension of the algebraic variety $\ker f = \{(\B{x},\B{y}) \in \overline{\FF}^{n_1} \times \overline{\FF}^{n_2} \mid f(\B{x},\B{y})=\B{0}\}$. 
\end{itemize}


We note that all three notions above generalize matrix rank. Moreover, just like matrix rank, for any $T \in \FF^{n \times n \times n}$ all three quantities lie in the range $[0,n]$. Furthermore, for the $n \times n \times n$ identity tensor $I_n$ we have $\SR(I_n)=\GR(I_n)=n$ and $\AR(I_n)=(1-o_{|\FF|}(1))n$.
The slice rank was defined by Tao~\cite{Tao16} in the context of the solution of the cap-set problem (similar notions have been considered before in other areas of research). 
The analytic rank was introduced by Gowers and Wolf~\cite{GowersWo11} in the context of higher-order Fourier analysis.
Roughly, this notion measures how close to uniform is the distribution of the values of the polynomial corresponding to the tensor.
The geometric rank is an algebro-geometric notion of rank that was recently introduced in~\cite{KoppartyMoZu20}. 
(A related notion was applied by Schmidt~\cite{Schmidt85} in the context of number theory.)
Intuitively, it measures the number of ``independent'' components of the corresponding bilinear map.
We use it as a geometric analogue 
of the bias of the 
(output distribution of the) 
bilinear map.


Understanding the structure of $d$-tensors, or $d$-dimensional matrices, that have low analytic rank is important in many applications of the structure-vs-randomness dichotomy, in additive combinatorics, coding theory and more (see, e.g.,~\cite{BhowmickLo15,GreenTao09}).
A recent breakthrough, obtained independently by Mili\'{c}evi\'{c}~\cite{Milicevic19} and by Janzer~\cite{Janzer19}, showed that the \emph{partition rank} of a $d$-tensor, which is a generalization of slice rank to $d$-tensors, is bounded from above by roughly $\AR(T)^{2^{2^{\poly(d)}}}$. For fixed $d$ this is a polynomial bound, 
which proves a conjecture of Kazhdan and Ziegler~\cite{KazhdanZi20}. 
Lovett~\cite{Lovett19}, as others have, asks whether in fact a linear upper bound holds.
For $3$-tensors, the best known bound until this work was $\SR(T) \le O(\AR(T)^4)$ by Haramaty and Shpilka~\cite{HaramatySh10}.
They write: ``It is an interesting open question to
decide whether we can do only with the $\sum_{j=1}^{O(\log_{|\FF|}1/\d)} \ell_i \cdot q_i$ part'', which refers to a linear upper bound $\SR(T) \le O(\AR(T))$.
Our main result is as follows.
\begin{main}[Main result]\label{main:summary}
	For any $3$-tensor $T$ over a field $\FF$, 
	$$\SR(T) \le 3\GR(T) \le 8.13\AR(T) $$
	where the first inequality holds over any perfect field\footnote{Commonly considered fields are perfect, including: any field of characteristic zero, any algebraically closed field, and any finite field.}, and the second for any finite field $\FF \neq \FF_2$.
\end{main}
